\uuid{8Ltr}
\titre{Trigonométrie}
\auteur{bodin ; hanani ; mzari}
\organisation{exo7}
\licence{CC BY-NC-SA 4.0}
\datecreate{2020-01-20}
\niveau{L1}
\module{Analyse}
\chapitre{Fonctions circulaires et hyperboliques inverses}
\sousChapitre{Autre}
\difficulte{3}
\theme{Fonctions usuelles, trigonométrie}
\exo7qcmid{lille-1.fonctions-usuelles.025}

\contenu{
\begin{question}
\qclassification{126.00}
\qtags{motcle=trigonométrie}
Soit $f(x)= \arcsin (\frac{2x}{1+x^2})$. On notera $D_f$ le domaine de définition de $f$. Quelles sont les assertions vraies ?
\begin{answers}
    \bad{$D_f=[-1,1]$}

    \good{$D_f=\Rr$}

    \good{Si $x \in [-1,1] $, $f(x) = 2 \arctan x$}

    \good{Si $x \ge 1 $, $f(x) = -2 \arctan x+\pi$}

\end{answers}
\begin{explanations}
$f$ est définie si $-1\le \frac{2x}{1+x^2} \le 1 $, c'est-à-dire $-1-x^2 \le 2x \le 1+x^2$, ce qui revient à : $(x+1)^2\ge 0$ et $(x-1)^2\ge 0$, ce qui est le cas pour tout $x \in \Rr$.

$f$ est dérivable sur $\Rr \backslash \{-1,1\}$ et
$f'(x) = \frac{2(1-x^2)}{|1-x^2|(1+x^2)} = \left\{\begin{array}{cc}\frac{2}{1+x^2},& \mbox{si} \, -1 < x < 1 \\ -\frac{2}{1+x^2},& \mbox{si} \, x > 1 \, \mbox{ou} \, x < -1 \end{array}\right. $. Comme $f$ est continue sur $\Rr$, on déduit que :
$f(x)= \left\{\begin{array}{ccc}2\arctan x + c_1, & \mbox{si} \, -1\le x\le 1 \\ -2\arctan x + c_2, & \mbox{si} \, x\le -1 \\ -2\arctan x + c_3, & \mbox{si} \, x\ge 1 \end{array}\right.$, où $c_1, c_2$ et $c_3$ sont des constantes.
En identifiant en $0$, en $-\infty$ et en $+\infty$, on obtient :
$f(x)= \left\{\begin{array}{ccc}2\arctan x, & \mbox{si} \, -1\le x\le 1 \\ -2\arctan x -\pi, & \mbox{si} \, x\le -1 \\ -2\arctan x +\pi, & \mbox{si} \, x\ge 1 \end{array}\right.$.
\end{explanations}
\end{question}
}
